\originalchapter{从经典波到物质波}\label{ch:quantum}
\originalsection{探测事件与概率振幅}
光强分布有干涉条纹, 但弱光探测仍以一个个离散事件出现. 电子逐个通过双缝时, 每次记录局域击中, 多次积累却形成波动条纹. 因而物质波的 $\psi$ 不是电子在每个位置的机械位移, 而是概率振幅; $|\psi|^2$ 在适当归一化后给位置概率密度. 对不能区分的两路径, 振幅相加 $\psi=\psi_1+\psi_2$, 概率含交叉项. 路径信息若与环境状态建立可辨相关, 交叉项可消失; 不能把所有干涉消失只解释成某个经典随机撞击.

先补上光强与单次探测之间的联系. 真空中频率为 $\nu$、波长为 $\lambda=c/\nu$ 的单个光子具有能量
\[
E_\gamma=h\nu=\frac{hc}{\lambda}=\hbar\omega,\qquad
\omega=2\pi\nu,\quad \hbar=\frac{h}{2\pi}.
\]
$h$ 是普朗克常量, $\hbar$ 是约化普朗克常量. 这是量子实验给出的能量关系, 不由经典波方程推出. 在频率不变时降低光强, 减少的是单位时间、单位面积的平均光子数 $I/E_\gamma$, 而非每个光子的能量; 增大波长才降低单光子能量, 同时使光学分辨率变差. 因此这两种“减弱照明”在路径探测中有不同后果.

设路径关联环境态 $|e_1\rangle,|e_2\rangle$. 符号 $|e\rangle$ (ket) 只表示一个归一化的量子状态向量; 可把环境状态想成一组振幅 $e_j$, 满足 $\sum_j|e_j|^2=1$. 内积 $\langle e_2|e_1\rangle=\sum_j e_{2j}^*e_{1j}$ 是两组振幅的复重叠, 模不超过1; 零重叠表示可以完全区分的状态. 连续变量时将求和换成积分即可. 总态为 $\psi_1(x)|e_1\rangle+\psi_2(x)|e_2\rangle$. 对未记录的环境逐项求概率并相加后, 位置概率为
\[|\psi_1|^2+|\psi_2|^2+2\Rea[\psi_1\psi_2^*\langle e_2|e_1\rangle].\] 环境态正交时, 即使不读出测量结果, 路径已可区分, 不出现该交叉项. 这一态空间说法的定义与运算见本仓库量子力学册第8章“态空间、测量与动力学”与第12章“复合系统、纠缠与全同粒子”. 本节只说明从光学干涉到物质波的接口, 不从经典波方程假称推出全部量子理论.

用辅助光照明电子路径是原课的具体思考实验. 增强照明可提高获得路径信息的事件比例, 条纹对比度随不可区分路径的比例下降; 降低光强则使多数电子未被散射, 在未获路径信息的子样本中仍可保留干涉. 若改用足够长波长, 衍射分辨率约为波长除以收集孔径, 两路径可能无法被光学分辨, 对应的环境态便有较大重叠. 这些比较说明“发生相互作用”和“留下可区分路径记录”需要分别判断, 不能仅凭照明存在便宣称条纹必然消失.

\originalsection{无限深势阱与量子色散}
质量 $m$ 的非相对论粒子在 $0<x<L$ 的无限深势阱中, 波函数边界为 $\psi(0)=\psi(L)=0$. 薛定谔动力学为
\[
\ii\hbar\psi_t=-\frac{\hbar^2}{2m}\psi_{xx}+V\psi.
\]
这是量子理论的动力学假设, 不由弦的二阶时间波方程推导. 阱内 $V=0$, 设 $\psi_n=A_n\sin(k_nx)\e^{-\ii\omega_nt}$, 得
\[
k_n=\frac{n\pi}{L},\quad A_n=\sqrt{2/L},\quad
\hbar\omega_n=E_n=\frac{\hbar^2k_n^2}{2m},\qquad n\ge1.
\]
空间模形状与两端固定弦相似, 频率却按 $n^2$ 而非 $n$ 增长. 相似边界不等于相同动力学. 概率密度 $|\psi_n|^2$ 在单一能量本征态中不随时间变化; 复相位旋转不是实弦点的上下简谐运动.

自由粒子色散为 $\omega=\hbar k^2/(2m)$, $p=\hbar k$, 因而群速度 $v_g=\hbar k/m=p/m$, 与经典粒子速度相接. 波包由许多 $k$ 的分量构成, 非零 $\omega''=\hbar/m$ 导致展宽. 无限阱定态的完整推导见量子力学册第5章“定态与一维势阱”. 空间--波数不确定关系给 $\Delta x\Delta p\ge\hbar/2$, 动量表象与算符证明见量子力学册第4章“动量表象与不确定性”. 原8.03SC第23讲文字材料的实质内容止于量子波与此桥梁; 课程页标题虽提引力波, 本轮未取得其对应文字推导, 不把未公开文字内容伪作已重编章节.
